\section{JUSTIFICATIVA}

A relevância desta pesquisa fundamenta-se em três pilares complementares: a demanda crescente por automação inteligente em operações de infraestrutura, a necessidade de preservação da segurança e privacidade dos dados, e a lacuna na literatura científica sobre a aplicação prática de LLMs na resolução de problemas de IaC.

Do ponto de vista operacional, relatórios da indústria indicam que equipes de DevOps gastam uma parcela significativa de seu tempo na resolução de incidentes e na análise de falhas de provisionamento \cite{puppet2023}. A capacidade de automatizar parcialmente esse processo, fornecendo diagnósticos iniciais e sugestões de correção, pode reduzir substancialmente o tempo médio de resolução (MTTR) e liberar engenheiros para atividades de maior valor agregado.

No aspecto de segurança, a execução local dos modelos de IA representa uma vantagem competitiva em relação a soluções baseadas em nuvem. Em setores regulamentados como o financeiro, de saúde e governamental, o envio de código de infraestrutura e logs operacionais para APIs externas pode violar normas como a LGPD (Lei Geral de Proteção de Dados) e políticas internas de segurança da informação. A abordagem proposta elimina esse risco ao manter todo o processamento no ambiente local.

Academicamente, embora existam estudos sobre a aplicação de LLMs na geração e depuração de código de aplicações, a literatura sobre seu uso específico na análise de falhas de IaC ainda é incipiente. Esta pesquisa contribui para preencher essa lacuna ao proporcionar evidências empíricas sobre a eficácia de diferentes modelos em cenários controlados, com metodologia reprodutível.
